\section{Introduction}
\label{sec:introduction}
Analogy supports reasoning by carrying a useful relational structure from a familiar source to a new target. In the classic fortress--radiation example, dividing an army across several routes suggests directing multiple low-intensity rays toward a tumor: individual paths remain safe, while their combined effect is sufficient at the target \citep{gick1980analogical}. The transferable content is a system of constraints and relations, including convergence and coordination. Sharing a topic or naming corresponding objects is insufficient.

Language models can generate examples, recognize analogies, and sometimes solve problems more accurately after receiving them. Analogical prompting uses self-generated problems and solutions \citep{yasunaga2024analogical}; teacher-generated analogies can also help scientific concept understanding \citep{yuan2024boosting}. However, random examples can produce comparable gains on some tasks \citep{qin2025relevant}. These findings raise a measurement question: \emph{relative to which controls, and through which operations, should an improvement count as evidence of analogical transfer?}

The question has two parts. First, the effect must be attributed carefully. Adding an analogy also adds text, explanation, and potentially answer-relevant knowledge. Comparing only against a prompt without examples cannot separate these factors. Even a positive average gain may combine repairs of incorrect answers with regressions on previously correct ones. Second, the process must be localized. A wrong answer can follow from inadequate sources, inconsistent correspondences, an unusable abstraction, or failure to apply a valid abstraction to target constraints. Final accuracy merges these failures. Conversely, a correct answer may be reached without using the analogy at all.

We address both parts through a theory-grounded framework for evaluating analogical reasoning in language models. Drawing on cognitive accounts of structural alignment, schema induction, and transfer, we operationalize four components: \textbf{retrieval}, \textbf{mapping}, \textbf{schema induction}, and \textbf{adaptation}. The model produces a separate output for each stage under an explicit information contract. We diagnose the first inadequate output in an unsuccessful trajectory, replace it with a reference-informed correction, and rerun the downstream stages with the same base model. This distinguishes an initial failure from additional failures that become visible after upstream repair.

The framework is supported by three connected studies. Phase~I examines the attribution of analogy effects in scientific question answering, including helpful and harmful answer transitions. Phase~II asks whether a solved, structurally corresponding source provides a selective advantage over a random source. Phase~III applies the stage framework to scientific and narrative problems to identify barriers to transfer. These phases answer different questions; together, they link controlled outcome measurement to process diagnosis.

Our contributions are:
\begin{itemize}
    \item A theory-grounded operational framework with explicit stage outputs and information boundaries, enabling diagnosis and cumulative intervention across retrieval, mapping, schema induction, and adaptation.
    \item A controlled analysis of analogy effects that distinguishes explanation, example length, relevance, textual mention, and item-level improvements and regressions.
    \item Empirical evidence for structure-dependent source-answer benefits, followed by a diagnosis showing that repairing early failures can reveal persistent downstream adaptation difficulties.
\end{itemize}
The stages are an externally imposed experimental decomposition. We neither claim to observe the model's hidden computation nor propose that every analogy must follow this temporal order. Our contribution is to make established theoretical distinctions testable in an auditable language-model workflow.
